\subsubsection{Subproceso 1: Acreditaci\'on y Pre-requisitos}

La Figura~\ref{fig:asis_1_flujo} muestra la secuencia horizontal de actividades y decisiones del proceso actual. El subproceso arranca d\'ias antes de la matr\'icula. Las Facultades cierran las actas de notas del ciclo anterior sin reclamos pendientes (RR-0570-2022, Art.~6\textdegree{}). En ese lapso, el estudiante cancela el autoseguro m\'edico obligatorio en Caja UNI o entidades bancarias (RR-2461-2023) y regulariza deudas con su Facultad (TUPA PA5400954D)~\cite{tupa2024,autoseguro2023}. En paralelo, la Escuela Profesional atiende dos situaciones especiales. Env\'ia a la DIRCE la n\'omina de alumnos egresantes con 30 cr\'editos o menos para otorgarles turno preferencial (Art.~42\textdegree{}). Los tutores, por su parte, autorizan el c\'odigo de quienes est\'an en riesgo acad\'emico tras cumplir su consejer\'ia previa (Arts.~8\textdegree{} y 48\textdegree{}).

Una vez cerradas las actas y recibidos los reportes, el sistema SIGA-ORCE cruza los pagos y valida los adeudos. Si el autoseguro no figura cancelado, el acceso queda bloqueado hasta que el estudiante subsane el pago. Con el requisito cumplido, el sistema calcula autom\'aticamente el promedio ponderado de los dos \'ultimos semestres regulares, fija el l\'imite m\'aximo de cr\'editos y establece el orden de m\'erito (Art.~9\textdegree{}, inc.~a)~\cite{reglamento2022}. Con estos datos, distribuye a los estudiantes en sus grupos de atenci\'on (Riesgo, Promoci\'on y Grupos Regulares) y genera la Ficha de Matr\'icula individual. El flujo culmina cuando la Intranet DIRCE publica el cronograma y muestra los cuatro controles de habilitaci\'on en verde, dejando la ficha lista para el inicio del Subproceso~2.

Los hallazgos del levantamiento del periodo 2026-2 quedan adem\'as incorporados al diagrama. Primero, el control de adeudos no oper\'o en el sistema: bast\'o pagar el autoseguro para quedar habilitado, sin bloquear por deudas de libros o equipos. Segundo, se identific\'o la falta de disponibilidad nocturna de los portales institucionales desde las 11:00~p.\,m. Este registro se tom\'o en periodo regular, fuera de la semana central de matr\'icula. Aun as\'i, durante la matr\'icula es frecuente que ocurran ca\'idas o saturaci\'on por alta concurrencia. Finalmente, los cambios de plan y convalidaciones se gestionan de forma f\'isica y manual ante la DIRCE, lo que genera retrasos para habilitar a los alumnos en su turno.

\begin{figure}[H]
    \centering
    \includegraphics[width=0.95\textwidth]{asis_subproceso_1_flujo.png}
    \caption{Diagrama de flujo (AS IS) del Subproceso 1}
    \label{fig:asis_1_flujo}
\end{figure}
